\begin{table}[t]
\centering
\small
\caption{Degree-bucket sanity check for fixed experts.}
\label{tab:degree_bucket_sanity_check}
\setlength{\tabcolsep}{3pt}
\begin{tabular}{p{0.10\textwidth}rp{0.19\textwidth}ccccp{0.25\textwidth}}
\toprule
Degree bucket & \#Queries & Dominant regime & Residual-only & Full Model & Gate-only & $\Delta$ Res.-Full & Interpretation \\
\midrule
0-1 & 98 & tail\_no\_img (86.7\%) & 0.0010 & 0.0787 & \textbf{0.0829} & -0.0777 & small sparse bucket; non-structural experts are stronger \\
2-5 & 938 & head\_has\_img (49.9\%) & 0.0017 & \textbf{0.0359} & 0.0283 & -0.0343 & low-degree mixed bucket favors the full model \\
6-20 & 9,433 & head\_has\_img (69.2\%) & 0.0025 & \textbf{0.0169} & 0.0164 & -0.0144 & head-dominated bucket; all experts remain weak \\
$>$20 & 9,531 & tail\_no\_img (98.3\%) & \textbf{0.6123} & 0.4130 & 0.3344 & +0.1993 & high-degree tail-side queries drive the structural advantage \\
\bottomrule
\end{tabular}
\vspace{0.4ex}
\caption*{\footnotesize \textit{Note:} Degree is the undirected count of training triples in which the target entity appears. MRR is averaged over three seed-specific evaluations; \#Queries reports the corresponding number of test queries per seed. The buckets expose the strong coupling between target degree and prediction-side regime: low- and mid-degree buckets are mostly head-side queries, whereas the highest-degree bucket is almost entirely tail-side.}
\end{table}
